\documentclass[10pt,letterpaper]{article}
\usepackage[utf8]{inputenc}
\usepackage[margin=0.6in]{geometry}
\usepackage{amsmath,amssymb,amsthm}
\usepackage{graphicx}
\usepackage{xcolor}
\usepackage{enumitem}
\usepackage{booktabs}
\usepackage{tabularx}
\usepackage{multicol}
\usepackage{tcolorbox}
\usepackage{fancyhdr}
\usepackage{titlesec}
\usepackage[hidelinks]{hyperref}
\usepackage{mdframed}

% Colors
\definecolor{criticalcolor}{RGB}{220,20,60}
\definecolor{examplecolor}{RGB}{0,100,0}
\definecolor{warningcolor}{RGB}{255,140,0}
\definecolor{headercolor}{RGB}{25,25,112}
\definecolor{formulacolor}{RGB}{0,51,102}

% Header and Footer
\pagestyle{fancy}
\fancyhf{}
\fancyhead[L]{\small\textbf{ORF309 Midterm 2 - THE ULTIMATE GUIDE}}
\fancyhead[R]{\small Page \thepage}
\fancyfoot[C]{\small Master these concepts, ace the exam!}

% Section formatting
\titleformat{\section}{\LARGE\bfseries\color{headercolor}}{\thesection}{1em}{}[\titlerule]
\titleformat{\subsection}{\Large\bfseries\color{headercolor}}{\thesubsection}{1em}{}
\titleformat{\subsubsection}{\large\bfseries\color{formulacolor}}{\thesubsubsection}{1em}{}

% Custom boxes
\newtcolorbox{criticalbox}[1][]{
  colback=criticalcolor!5,
  colframe=criticalcolor,
  fonttitle=\bfseries,
  title={\textcolor{white}{★ CRITICAL - MASTER THIS FIRST}},
  #1
}

\newtcolorbox{examplebox}[1][]{
  colback=examplecolor!5,
  colframe=examplecolor!80,
  fonttitle=\bfseries,
  title={WORKED EXAMPLE},
  #1
}

\newtcolorbox{formulabox}[1][]{
  colback=formulacolor!5,
  colframe=formulacolor,
  fonttitle=\bfseries,
  title={KEY FORMULAS},
  #1
}

\newtcolorbox{warningbox}[1][]{
  colback=warningcolor!5,
  colframe=warningcolor,
  fonttitle=\bfseries,
  title={⚠ COMMON MISTAKE - AVOID THIS!},
  #1
}

\newtcolorbox{strategybox}[1][]{
  colback=blue!5,
  colframe=blue!70,
  fonttitle=\bfseries,
  title={PROBLEM-SOLVING STRATEGY},
  #1
}

% Compact lists
\setlist{nosep, leftmargin=*}

\begin{document}

% Title Page
\begin{titlepage}
\centering
\vspace*{1cm}

{\Huge\bfseries\color{criticalcolor} THE ULTIMATE\par}
\vspace{0.3cm}
{\Huge\bfseries ORF309 MIDTERM 2\par}
\vspace{0.3cm}
{\Huge\bfseries STUDY GUIDE\par}

\vspace{1cm}
{\Large\textit{Everything You Need to Ace This Exam}\par}

\vspace{1.5cm}

\begin{tcolorbox}[colback=yellow!10,colframe=orange,width=0.9\textwidth,arc=3mm]
\centering
\large
\textbf{One Guide. All Topics. Step-by-Step.}\\[0.3cm]
From Zero to Hero in 3-4 Days
\end{tcolorbox}

\vspace{1cm}

{\large\textbf{What's Inside:}\par}
\begin{itemize}[leftmargin=2cm]
  \item ✓ Beginner-friendly explanations with real-world examples
  \item ✓ Every formula you need (color-coded by importance)
  \item ✓ 15+ fully worked examples from past exams
  \item ✓ Pattern recognition guide (when to use what technique)
  \item ✓ Common mistakes and how to avoid them
  \item ✓ Step-by-step problem-solving strategies
\end{itemize}

\vfill

{\large Created: March 31, 2026\par}
{\large Based on analysis of 7 past Midterm 2 exams\par}

\end{titlepage}

\tableofcontents
\newpage

% Quick Start Guide
\section*{START HERE: Quick Navigation}
\addcontentsline{toc}{section}{START HERE: Quick Navigation}

\begin{tcolorbox}[colback=yellow!20,colframe=orange!80,arc=3mm]
\textbf{First Time Learning?} Read in order: Section 1 → 2 → 3 → 4

\textbf{Need Quick Review?} Go straight to ``Formula Sheets'' (Section 8)

\textbf{Doing Practice Problems?} Use ``Pattern Recognition Guide'' (Section 7)

\textbf{Last-Minute Cramming?} Focus on Sections 1, 2, and 8 only
\end{tcolorbox}

\subsection*{The Big 3 Topics (Guaranteed on Exam)}

\begin{enumerate}[leftmargin=*]
  \item \textbf{Poisson Processes} (Section 1) - 40\% of exam, appears in 100\% of exams
    \begin{itemize}
      \item Master thinning, competing processes, conditional expectations
    \end{itemize}

  \item \textbf{Random Walks} (Section 2) - 35\% of exam, appears in 100\% of exams
    \begin{itemize}
      \item Master first-step analysis for hitting probabilities and expected times
    \end{itemize}

  \item \textbf{Joint/Conditional Densities} (Section 3) - 25\% of exam, appears in 86\% of exams
    \begin{itemize}
      \item Master law of total expectation and building joint PDFs
    \end{itemize}
\end{enumerate}

\subsection*{How to Use This Guide}

\textbf{Day 1:} Read Section 1 (Poisson), do all examples, try 2 practice problems

\textbf{Day 2:} Read Section 2 (Random Walks), do all examples, try 2 practice problems

\textbf{Day 3:} Read Section 3 (Joint Densities), review Sections 4-5, do full practice exam

\textbf{Day 4:} Review Section 8 (formulas), do another practice exam, make cheat sheet

\newpage

\section{POISSON PROCESSES - The Most Important Topic}

\begin{criticalbox}
\textbf{Why This Matters:} Appears in \textbf{7/7 exams (100\%)}. Average weight: \textbf{40\% of total points}.

Master thinning and conditional expectations and you'll solve half the exam!
\end{criticalbox}

\subsection{What is a Poisson Process?}

\subsubsection{The Big Idea (In Plain English)}

Imagine sitting at Starbucks counting customers. Sometimes 5 people arrive in 2 minutes, sometimes nobody for 10 minutes. Arrivals seem random, but there's an \textit{average rate}.

A \textbf{Poisson process} models these random arrivals:
\begin{itemize}
  \item Events happen randomly over time
  \item There's an average rate $\lambda$ (like ``3 customers per minute'')
  \item Each event is independent of others
\end{itemize}

\subsubsection{The Basic Math}

If events arrive at rate $\lambda$ (per unit time):

\begin{formulabox}
\textbf{Number of arrivals in time $t$:} $N(t) \sim \text{Poisson}(\lambda t)$

\vspace{0.2cm}

\textbf{Probability of exactly $k$ arrivals:}
\[
P(N(t) = k) = \frac{(\lambda t)^k e^{-\lambda t}}{k!}
\]

\textbf{Expected number of arrivals:} $E[N(t)] = \lambda t$

\textbf{Variance:} $\text{Var}(N(t)) = \lambda t$
\end{formulabox}

\begin{examplebox}
\textbf{Problem:} Emails arrive at your inbox at rate $\lambda = 10$ per hour. How many emails do you expect in 3 hours?

\textbf{Solution:}
\[
E[N(3)] = \lambda t = 10 \times 3 = \boxed{30 \text{ emails}}
\]

\textbf{What's the probability of exactly 25 emails in 3 hours?}
\begin{align*}
P(N(3) = 25) &= \frac{(10 \cdot 3)^{25} e^{-10 \cdot 3}}{25!} \\
&= \frac{30^{25} e^{-30}}{25!} \approx 0.0387
\end{align*}
\end{examplebox}

\subsection{THINNING - The Most Tested Concept}

\begin{criticalbox}
\textbf{Appears in 6/7 exams!} If you master ONE thing, make it thinning.
\end{criticalbox}

\subsubsection{The Setup}

Events arrive at rate $\lambda$. Each event is independently classified as:
\begin{itemize}
  \item ``Type A'' with probability $p$
  \item ``Type B'' with probability $1-p$
\end{itemize}

\textbf{Question:} What's the rate of Type A events?

\subsubsection{The Answer (Thinning Property)}

\begin{formulabox}[colback=yellow!20,colframe=orange!80]
\textbf{THINNING PROPERTY:}

If arrivals are Poisson($\lambda$) and each is Type A with probability $p$:
\[
\boxed{\text{Type A arrivals are Poisson}(\lambda p)}
\]

You ``thin'' the arrival rate by multiplying by the probability!
\end{formulabox}

\subsubsection{Intuition}

If 10 customers arrive per minute and 30\% order coffee:
\[
\text{Coffee orders} = 10 \times 0.3 = 3 \text{ per minute}
\]

\begin{examplebox}
\textbf{Problem (Exam Style):} At a coffee shop, customers arrive at rate 2 per minute. Each customer:
\begin{itemize}
  \item Orders coffee with probability 0.4
  \item Orders tea with probability 0.4
  \item Orders food with probability 0.2
\end{itemize}

\textbf{(a)} What's the rate of coffee orders?

\textbf{Solution (a):} Use thinning:
\[
\lambda_{\text{coffee}} = 2 \times 0.4 = \boxed{0.8 \text{ per minute}}
\]

\textbf{(b)} On average, how many coffee orders in 10 minutes?

\textbf{Solution (b):}
\[
E[\text{coffee in 10 min}] = \lambda_{\text{coffee}} \times t = 0.8 \times 10 = \boxed{8 \text{ coffee orders}}
\]

\textbf{(c)} What's the probability that exactly 4 coffee orders happen before the first food order?

\textbf{Solution (c):} Think of it as a race:
\begin{itemize}
  \item Coffee orders: Poisson(0.8/min)
  \item Food orders: Poisson(0.4/min)
\end{itemize}

For exactly 4 coffee before first food, we want the sequence: Coffee, Coffee, Coffee, Coffee, Food.

\textbf{Key insight:} Tea orders don't matter! Among coffee+food orders only, each one is independently:
\begin{itemize}
  \item Coffee with probability: $\frac{0.8}{0.8+0.4} = \frac{2}{3}$
  \item Food with probability: $\frac{0.4}{0.8+0.4} = \frac{1}{3}$
\end{itemize}

So we need the specific sequence C-C-C-C-F:
\[
P(\text{4 coffee before 1 food}) = \left(\frac{2}{3}\right)^4 \left(\frac{1}{3}\right) = \frac{16}{81 \times 3} = \boxed{\frac{16}{243}}
\]

\textbf{Why doesn't tea matter?} Tea customers can arrive in between, but we only care about the coffee vs food "race." It's like watching only red and blue cars pass by and asking "what's the probability I see 4 red before 1 blue?" - green cars (tea) passing by don't affect this.
\end{examplebox}

\subsection{Competing Processes}

\subsubsection{The Setup}

Two independent Poisson processes:
\begin{itemize}
  \item Process A at rate $\lambda_A$
  \item Process B at rate $\lambda_B$
\end{itemize}

\textbf{Question:} Which happens first?

\subsubsection{The Formulas}

\begin{formulabox}
\textbf{COMPETING PROCESSES:}

\vspace{0.2cm}

\textbf{Probability A happens before B:}
\[
P(\text{A before B}) = \frac{\lambda_A}{\lambda_A + \lambda_B}
\]

\textbf{Expected time until first event (A or B):}
\[
E[T] = \frac{1}{\lambda_A + \lambda_B}
\]

\textbf{The first event to occur:} Poisson($\lambda_A + \lambda_B$)
\end{formulabox}

\begin{examplebox}
\textbf{Problem:} You're waiting for:
\begin{itemize}
  \item Uber: arrives at rate $\lambda_U = 6$ per hour
  \item Bus: arrives at rate $\lambda_B = 4$ per hour
\end{itemize}

\textbf{(a)} What's the probability Uber arrives first?

\textbf{Solution (a):}
\[
P(\text{Uber first}) = \frac{6}{6+4} = \frac{6}{10} = \boxed{0.6 = 60\%}
\]

\textbf{(b)} How long do you expect to wait?

\textbf{Solution (b):}
\[
E[T] = \frac{1}{\lambda_U + \lambda_B} = \frac{1}{6+4} = \frac{1}{10} \text{ hour} = \boxed{6 \text{ minutes}}
\]
\end{examplebox}

\subsection{Conditional Expectations with Poisson}

\begin{criticalbox}
\textbf{Appears in 5/7 exams!} This is almost as important as thinning.
\end{criticalbox}

\subsubsection{The Setup}

Events arrive at rate $\lambda$. You're told that exactly $n$ events occurred.

\textbf{Question:} What's the expected time?

\subsubsection{The Formula}

\begin{formulabox}[colback=yellow!20,colframe=orange!80]
\textbf{CONDITIONAL EXPECTATION:}

Given that exactly $n$ arrivals occurred by time $T$:
\[
\boxed{E[T \mid N(T) = n] = \frac{n}{\lambda}}
\]

\textbf{Intuition:} If $n$ events happened at rate $\lambda$, it must have taken $n/\lambda$ time.
\end{formulabox}

\begin{examplebox}
\textbf{Problem:} Customers arrive at rate $\lambda = 2$ per minute. Given that exactly 10 customers arrived, what's the expected time?

\textbf{Solution:}
\[
E[T \mid N(T) = 10] = \frac{10}{2} = \boxed{5 \text{ minutes}}
\]

\textbf{Problem (Harder):} At a taxi stop, taxis arrive at rate 1/min, passengers at rate 2/min. Given that by the time the first taxi arrives, exactly 10 passengers are waiting, what's the expected time?

\textbf{Solution:}
\begin{itemize}
  \item First taxi arrival time: $T \sim \text{Exp}(1)$
  \item Given 10 passengers arrived: $E[T \mid N_{\text{pass}}(T) = 10] = \frac{10}{2} = \boxed{5 \text{ minutes}}$
\end{itemize}
\end{examplebox}

\subsection{Poisson Quick Reference}

\begin{formulabox}
\textbf{POISSON PROCESSES - FORMULA SHEET}

\vspace{0.2cm}

\textbf{Basic Facts:}
\begin{align*}
E[N(t)] &= \lambda t \\
\text{Var}(N(t)) &= \lambda t \\
P(N(t) = k) &= \frac{(\lambda t)^k e^{-\lambda t}}{k!}
\end{align*}

\textbf{Thinning:} Type A prob $p$ → Type A is Poisson($\lambda p$)

\textbf{Competing:} $P(\text{A before B}) = \dfrac{\lambda_A}{\lambda_A + \lambda_B}$

\textbf{Conditional:} $E[T \mid N(T) = n] = \dfrac{n}{\lambda}$
\end{formulabox}

\begin{warningbox}
\textbf{Common Mistakes:}
\begin{enumerate}
  \item Forgetting to use thinning when events split into types
  \item Using $\lambda$ instead of $\lambda t$ in Poisson formula
  \item Not recognizing conditional expectation problems
\end{enumerate}
\end{warningbox}

\newpage

\section{RANDOM WALKS \& FIRST-STEP ANALYSIS}

\begin{criticalbox}
\textbf{Why This Matters:} Appears in \textbf{7/7 exams (100\%)}. Average weight: \textbf{35\% of total points}.

Master first-step analysis and you've conquered another third of the exam!
\end{criticalbox}

\subsection{What is a Random Walk?}

\subsubsection{The Big Idea (In Plain English)}

Imagine a gambling game:
\begin{itemize}
  \item You start with \$50
  \item Each round: flip a coin
  \item Heads: win \$1, Tails: lose \$1
  \item Game ends at \$0 (bankrupt) or \$100 (rich)
\end{itemize}

Your money randomly walks up and down. This is a \textbf{random walk}!

\subsubsection{Formal Setup}

\begin{itemize}
  \item \textbf{States:} Positions $a, a+1, a+2, \ldots, b-1, b$
  \item \textbf{Barriers:} $a$ (lower) and $b$ (upper) where game ends
  \item \textbf{Starting position:} $S_0 = i$ (somewhere between $a$ and $b$)
  \item \textbf{Transitions:} Each step, move up/down/stay with given probabilities
\end{itemize}

\subsection{The Two Main Questions}

Every random walk problem asks one (or both) of:

\begin{enumerate}
  \item \textbf{Hitting Probability:} What's $P(\text{reach } b \text{ before } a)$?
  \item \textbf{Expected Hitting Time:} How long until we reach $a$ or $b$?
\end{enumerate}

\subsection{FIRST-STEP ANALYSIS - The Key Technique}

\begin{criticalbox}
This technique appears in \textbf{100\% of exams}. Master it and random walk problems become easy!
\end{criticalbox}

\subsubsection{The Core Idea}

To figure out what happens from position $i$:
\begin{enumerate}
  \item Think about what happens on the \textit{first step}
  \item Use what you know about where you'll be next
  \item Write a recursive equation
\end{enumerate}

It's like saying: ``My chance of winning from \$50 depends on whether I go to \$51 or \$49 next, and my chances from those positions.''

\subsection{First-Step Analysis: HITTING PROBABILITY}

\subsubsection{Setup}

Let $p(i)$ = probability of reaching $b$ before $a$, starting from position $i$.

\subsubsection{The Recursive Formula}

\begin{formulabox}[colback=yellow!20,colframe=orange!80]
\textbf{FIRST-STEP ANALYSIS FOR HITTING PROBABILITY:}

\vspace{0.2cm}

\textbf{Recursive formula:}
\[
\boxed{p(i) = p_{\uparrow} \cdot p(i+1) + p_{\downarrow} \cdot p(i-1) + p_{\leftrightarrow} \cdot p(i)}
\]

where $p_{\uparrow}$ = prob of moving up, $p_{\downarrow}$ = prob of moving down, $p_{\leftrightarrow}$ = prob of staying

\vspace{0.2cm}

\textbf{Boundary conditions:}
\[
\boxed{p(a) = 0 \quad \text{and} \quad p(b) = 1}
\]

\textbf{Solve the system of equations for all intermediate states!}
\end{formulabox}

\subsubsection{Special Case: Symmetric Random Walk}

When $p_{\uparrow} = p_{\downarrow} = 1/2$ (fair coin flip):

\begin{formulabox}
\textbf{SYMMETRIC WALK SHORTCUT:}
\[
\boxed{p(i) = \frac{i - a}{b - a}}
\]

Your hitting probability is just your relative position between the barriers!
\end{formulabox}

\begin{examplebox}
\textbf{Problem:} You play a game with positions 0, 1, 2, 3, 4. Each round:
\begin{itemize}
  \item Move up with probability $1/2$
  \item Move down with probability $1/2$
\end{itemize}

Start at position 2. Win at 4, lose at 0.

\textbf{What's the probability of winning?}

\textbf{Solution (Method 1 - Shortcut):}

This is symmetric, so:
\[
p(2) = \frac{2 - 0}{4 - 0} = \frac{2}{4} = \boxed{0.5 = 50\%}
\]

\textbf{Solution (Method 2 - First-Step Analysis):}

Boundaries: $p(0) = 0$, $p(4) = 1$

From position 1:
\[
p(1) = \frac{1}{2} p(2) + \frac{1}{2} p(0) = \frac{1}{2} p(2)
\]

From position 2:
\[
p(2) = \frac{1}{2} p(3) + \frac{1}{2} p(1)
\]

From position 3:
\[
p(3) = \frac{1}{2} p(4) + \frac{1}{2} p(2) = \frac{1}{2} + \frac{1}{2} p(2)
\]

Substitute $p(1) = \frac{1}{2}p(2)$ and $p(3) = \frac{1}{2} + \frac{1}{2}p(2)$ into $p(2)$ equation:
\begin{align*}
p(2) &= \frac{1}{2}\left(\frac{1}{2} + \frac{1}{2}p(2)\right) + \frac{1}{2}\left(\frac{1}{2}p(2)\right) \\
p(2) &= \frac{1}{4} + \frac{1}{4}p(2) + \frac{1}{4}p(2) \\
p(2) &= \frac{1}{4} + \frac{1}{2}p(2) \\
\frac{1}{2}p(2) &= \frac{1}{4} \\
p(2) &= \boxed{0.5}
\end{align*}
\end{examplebox}

\subsection{First-Step Analysis: EXPECTED TIME}

\begin{criticalbox}
\textbf{Most Common Mistake:} Forgetting the ``+1'' in this formula! Don't be that person!
\end{criticalbox}

\subsubsection{Setup}

Let $\mu(i)$ = expected number of steps to reach $a$ or $b$, starting from position $i$.

\subsubsection{The Recursive Formula}

\begin{formulabox}[colback=red!10,colframe=red!80]
\textbf{FIRST-STEP ANALYSIS FOR EXPECTED TIME:}

\vspace{0.2cm}

\textbf{Recursive formula:}
\[
\boxed{\mu(i) = \mathbf{1} + p_{\uparrow} \cdot \mu(i+1) + p_{\downarrow} \cdot \mu(i-1) + p_{\leftrightarrow} \cdot \mu(i)}
\]

\textbf{DON'T FORGET THE ``+1''!} (Counts the current step)

\vspace{0.2cm}

\textbf{Boundary conditions:}
\[
\boxed{\mu(a) = 0 \quad \text{and} \quad \mu(b) = 0}
\]

(Game is over at the barriers!)
\end{formulabox}

\subsubsection{Special Case: Symmetric Random Walk}

When $p_{\uparrow} = p_{\downarrow} = 1/2$:

\begin{formulabox}
\textbf{SYMMETRIC WALK SHORTCUT:}
\[
\boxed{\mu(i) = (i - a)(b - i)}
\]
\end{formulabox}

\begin{examplebox}
\textbf{Problem:} Same game as before (positions 0--4, symmetric). Start at 2.

\textbf{How many steps on average until the game ends?}

\textbf{Solution (Method 1 - Shortcut):}
\[
\mu(2) = (2 - 0)(4 - 2) = 2 \times 2 = \boxed{4 \text{ steps}}
\]

\textbf{Solution (Method 2 - First-Step Analysis):}

Boundaries: $\mu(0) = 0$, $\mu(4) = 0$

From position 1:
\[
\mu(1) = 1 + \frac{1}{2}\mu(2) + \frac{1}{2}\mu(0) = 1 + \frac{1}{2}\mu(2)
\]

From position 2:
\[
\mu(2) = 1 + \frac{1}{2}\mu(3) + \frac{1}{2}\mu(1)
\]

From position 3:
\[
\mu(3) = 1 + \frac{1}{2}\mu(4) + \frac{1}{2}\mu(2) = 1 + \frac{1}{2}\mu(2)
\]

Notice $\mu(1) = \mu(3) = 1 + \frac{1}{2}\mu(2)$ by symmetry.

Substitute into $\mu(2)$ equation:
\begin{align*}
\mu(2) &= 1 + \frac{1}{2}(1 + \frac{1}{2}\mu(2)) + \frac{1}{2}(1 + \frac{1}{2}\mu(2)) \\
\mu(2) &= 1 + \frac{1}{2} + \frac{1}{4}\mu(2) + \frac{1}{2} + \frac{1}{4}\mu(2) \\
\mu(2) &= 2 + \frac{1}{2}\mu(2) \\
\frac{1}{2}\mu(2) &= 2 \\
\mu(2) &= \boxed{4}
\end{align*}
\end{examplebox}

\subsection{Biased Random Walks}

\subsubsection{When Probabilities Are Not Equal}

Sometimes $P(\text{up}) \neq P(\text{down})$. This is a \textbf{biased random walk}.

\textbf{Approach:} Use first-step analysis and solve the system of equations (no shortcuts!).

\begin{examplebox}
\textbf{Problem (Exam-Style):} Positions 0, 1, 2, 3. Each step:
\begin{itemize}
  \item Move up with probability $2/3$
  \item Move down with probability $1/3$
\end{itemize}

Start at position 1. What's $p(1)$?

\textbf{Solution:}

Boundaries: $p(0) = 0$, $p(3) = 1$

From position 1:
\[
p(1) = \frac{2}{3}p(2) + \frac{1}{3}p(0) = \frac{2}{3}p(2)
\]

From position 2:
\[
p(2) = \frac{2}{3}p(3) + \frac{1}{3}p(1) = \frac{2}{3} + \frac{1}{3}p(1)
\]

Substitute $p(1) = \frac{2}{3}p(2)$ into second equation:
\begin{align*}
p(2) &= \frac{2}{3} + \frac{1}{3} \cdot \frac{2}{3}p(2) \\
p(2) &= \frac{2}{3} + \frac{2}{9}p(2) \\
\frac{7}{9}p(2) &= \frac{2}{3} \\
p(2) &= \frac{2}{3} \cdot \frac{9}{7} = \frac{6}{7}
\end{align*}

Therefore:
\[
p(1) = \frac{2}{3} \cdot \frac{6}{7} = \boxed{\frac{4}{7}}
\]
\end{examplebox}

\subsection{Random Walks Quick Reference}

\begin{formulabox}
\textbf{RANDOM WALKS - FORMULA SHEET}

\vspace{0.2cm}

\textbf{Hitting Probability $p(i)$:}
\[
p(i) = p_{\uparrow} p(i+1) + p_{\downarrow} p(i-1) + p_{\leftrightarrow} p(i)
\]
Boundaries: $p(a) = 0$, $p(b) = 1$

\textbf{Expected Time $\mu(i)$:}
\[
\mu(i) = \mathbf{1} + p_{\uparrow} \mu(i+1) + p_{\downarrow} \mu(i-1) + p_{\leftrightarrow} \mu(i)
\]
Boundaries: $\mu(a) = 0$, $\mu(b) = 0$

\vspace{0.2cm}

\textbf{Symmetric Walk Shortcuts ($p_{\uparrow} = p_{\downarrow} = 1/2$):}
\begin{align*}
p(i) &= \frac{i-a}{b-a} \\
\mu(i) &= (i-a)(b-i)
\end{align*}
\end{formulabox}

\begin{warningbox}
\textbf{Common Mistakes:}
\begin{enumerate}
  \item \textbf{Forgetting ``+1''} in $\mu(i)$ formula (MOST COMMON!)
  \item Wrong boundary conditions ($\mu$ should be 0 at BOTH barriers)
  \item Not setting up all equations (need one for each intermediate state)
  \item Arithmetic errors when solving systems
\end{enumerate}
\end{warningbox}

\begin{strategybox}
\textbf{HOW TO SOLVE RANDOM WALK PROBLEMS:}

\textbf{Step 1:} Draw a diagram with all states

\textbf{Step 2:} Identify what you're solving for ($p(i)$ or $\mu(i)$?)

\textbf{Step 3:} Check if symmetric → use shortcut if possible

\textbf{Step 4:} Write boundary conditions

\textbf{Step 5:} Write first-step equation for each intermediate state

\textbf{Step 6:} Solve the system (usually 2-3 equations, 2-3 unknowns)

\textbf{Step 7:} Check answer makes intuitive sense
\end{strategybox}

\newpage

\section{JOINT \& CONDITIONAL DENSITIES}

\begin{criticalbox}
\textbf{Why This Matters:} Appears in \textbf{6/7 exams (86\%)}. Average weight: \textbf{25\% of total points}.
\end{criticalbox}

\subsection{What Are We Dealing With?}

\subsubsection{The Big Idea}

Sometimes we have \textbf{two related random variables}:
\begin{itemize}
  \item $X$ = total restaurant bill
  \item $Y$ = amount you pay (depends on $X$!)
\end{itemize}

We need to understand how they relate using \textbf{joint} and \textbf{conditional} densities.

\subsection{Key Definitions}

\begin{formulabox}
\textbf{JOINT PDF:} $f_{X,Y}(x,y)$ describes both variables together

\textbf{MARGINAL PDF:} Distribution of one variable alone
\begin{align*}
f_X(x) &= \int_{-\infty}^{\infty} f_{X,Y}(x,y) \, dy \\
f_Y(y) &= \int_{-\infty}^{\infty} f_{X,Y}(x,y) \, dx
\end{align*}

\textbf{CONDITIONAL PDF:} Distribution of one given the other
\[
f_{X|Y}(x|y) = \frac{f_{X,Y}(x,y)}{f_Y(y)}, \quad f_{Y|X}(y|x) = \frac{f_{X,Y}(x,y)}{f_X(x)}
\]

\textbf{BUILDING JOINT FROM CONDITIONAL:}
\[
f_{X,Y}(x,y) = f_{Y|X}(y|x) \cdot f_X(x) = f_{X|Y}(x|y) \cdot f_Y(y)
\]
\end{formulabox}

\subsection{Law of Total Expectation}

\begin{criticalbox}
This is one of the most powerful tools in probability. Master it!
\end{criticalbox}

\begin{formulabox}[colback=yellow!20,colframe=orange!80]
\textbf{LAW OF TOTAL EXPECTATION (Tower Property):}
\[
\boxed{E[Y] = E[E[Y|X]]}
\]

In integral form:
\[
E[Y] = \int_{-\infty}^{\infty} E[Y|X=x] \cdot f_X(x) \, dx
\]

\textbf{When to use:} When $Y$ depends on $X$ and you know the conditional expectation.
\end{formulabox}

\subsection{The $Y = UX$ Transformation (Super Common!)}

This appears in many exam problems:

\begin{itemize}
  \item $X$ = some random variable (e.g., bill amount, loss amount)
  \item $U \sim \text{Uniform}[0,1]$, independent of $X$
  \item $Y = UX$ = fraction of $X$
\end{itemize}

\begin{formulabox}
\textbf{KEY FACTS FOR $Y = UX$:}

\vspace{0.2cm}

Given $X = x$:
\begin{itemize}
  \item $Y = Ux$ is uniform on $[0, x]$
  \item $f_{Y|X}(y|x) = \frac{1}{x}$ for $0 < y < x$
  \item $E[Y|X=x] = \frac{x}{2}$
\end{itemize}

\textbf{Joint PDF:}
\[
f_{X,Y}(x,y) = f_{Y|X}(y|x) \cdot f_X(x) = \frac{1}{x} f_X(x) \quad \text{for } 0 < y < x
\]

\textbf{Always specify the range!}
\end{formulabox}

\begin{examplebox}
\textbf{Problem (Exam-Style):}
\begin{itemize}
  \item $X$ = total loss (in \$1000s), with $P(X > x) = 1/x^2$ for $x \geq 1$
  \item $U \sim \text{Uniform}[0,1]$, independent of $X$
  \item $Y = UX$ = amount you pay out-of-pocket
\end{itemize}

\textbf{(a)} Find $E[Y]$ (expected payment).

\textbf{Solution (a):} Use law of total expectation!

Step 1: Find $E[Y|X=x]$
\[
E[Y|X=x] = E[Ux|X=x] = xE[U] = x \cdot \frac{1}{2} = \frac{x}{2}
\]

Step 2: Find $f_X(x)$
\[
P(X > x) = 1/x^2 \Rightarrow F_X(x) = 1 - 1/x^2
\]
\[
f_X(x) = \frac{d}{dx}F_X(x) = \frac{2}{x^3} \quad \text{for } x \geq 1
\]

Step 3: Apply law of total expectation
\begin{align*}
E[Y] &= E[E[Y|X]] = E\left[\frac{X}{2}\right] = \frac{1}{2}E[X] \\
E[X] &= \int_1^{\infty} x \cdot \frac{2}{x^3} \, dx = \int_1^{\infty} \frac{2}{x^2} \, dx \\
&= \left[-\frac{2}{x}\right]_1^{\infty} = 0 - (-2) = 2
\end{align*}

Therefore:
\[
E[Y] = \frac{1}{2} \cdot 2 = \boxed{\$1000}
\]

\textbf{(b)} Find joint PDF $f_{X,Y}(x,y)$.

\textbf{Solution (b):}
\[
f_{X,Y}(x,y) = f_{Y|X}(y|x) \cdot f_X(x) = \frac{1}{x} \cdot \frac{2}{x^3} = \boxed{\frac{2}{x^4}} \quad \text{for } 0 < y < x, \, x \geq 1
\]

\textbf{CRITICAL: Always specify the range!}

\textbf{(c)} Given you paid exactly \$1000 ($Y=1$), what's $P(X > 2)$?

\textbf{Solution (c):} Need conditional probability $P(X > 2 | Y = 1)$.

First find $f_Y(1)$:
\begin{align*}
f_Y(1) &= \int_1^{\infty} f_{X,Y}(1, y) \, dx \\
&= \int_1^{\infty} \frac{2}{x^4} \, dx = \left[-\frac{2}{3x^3}\right]_1^{\infty} = \frac{2}{3}
\end{align*}

Then:
\begin{align*}
P(X > 2 | Y = 1) &= \frac{\int_2^{\infty} f_{X,Y}(x, 1) \, dx}{f_Y(1)} \\
&= \frac{\int_2^{\infty} \frac{2}{x^4} \, dx}{2/3} = \frac{[-2/(3x^3)]_2^{\infty}}{2/3} \\
&= \frac{2/(3 \cdot 8)}{2/3} = \frac{1/12}{2/3} = \boxed{\frac{1}{8}}
\end{align*}
\end{examplebox}

\subsection{Joint Densities Quick Reference}

\begin{formulabox}
\textbf{JOINT/CONDITIONAL DENSITIES - FORMULA SHEET}

\vspace{0.2cm}

\textbf{Relationships:}
\begin{align*}
f_{X|Y}(x|y) &= \frac{f_{X,Y}(x,y)}{f_Y(y)} \\
f_{X,Y}(x,y) &= f_{Y|X}(y|x) \cdot f_X(x) \\
f_X(x) &= \int f_{X,Y}(x,y) \, dy
\end{align*}

\textbf{Law of Total Expectation:}
\[
E[Y] = E[E[Y|X]]
\]

\textbf{For $Y = UX$ where $U \sim \text{Unif}[0,1]$:}
\begin{align*}
f_{Y|X}(y|x) &= \frac{1}{x} \quad (0 < y < x) \\
E[Y|X=x] &= \frac{x}{2} \\
f_{X,Y}(x,y) &= \frac{f_X(x)}{x} \quad (0 < y < x)
\end{align*}
\end{formulabox}

\begin{warningbox}
\textbf{Common Mistakes:}
\begin{enumerate}
  \item Not specifying the range where $f_{X,Y}(x,y) > 0$
  \item Forgetting to find marginal $f_Y(y)$ when computing conditional
  \item Sign errors when computing CDFs from tail probabilities
  \item Not checking that PDF integrates to 1
\end{enumerate}
\end{warningbox}

\begin{strategybox}
\textbf{PROBLEM-SOLVING STRATEGY:}

\textbf{Step 1:} Identify what's given (marginal? conditional?)

\textbf{Step 2:} If computing expectation, try law of total expectation first

\textbf{Step 3:} If building joint PDF, use $f_{X,Y} = f_{Y|X} \cdot f_X$

\textbf{Step 4:} Always specify range! Draw picture if needed

\textbf{Step 5:} Sanity check: Does it integrate to 1? Makes sense?
\end{strategybox}

\newpage

\section{EXPONENTIAL DISTRIBUTION \& LIFETIMES}

\begin{criticalbox}
\textbf{Frequency:} Appears in \textbf{5/7 exams (71\%)}. Often combined with Poisson.
\end{criticalbox}

\subsection{What is Exponential Distribution?}

\subsubsection{The Big Idea}

Exponential distribution models \textbf{waiting times}:
\begin{itemize}
  \item Time until next customer arrives
  \item Lifetime of a light bulb
  \item Time until your phone buzzes
\end{itemize}

\textbf{Key parameter:} $\lambda$ = rate (events per unit time)

\subsection{Key Facts}

\begin{formulabox}
\textbf{EXPONENTIAL DISTRIBUTION $T \sim \text{Exp}(\lambda)$:}

\vspace{0.2cm}

\textbf{PDF:} $f(t) = \lambda e^{-\lambda t}$ for $t > 0$

\textbf{CDF:} $F(t) = 1 - e^{-\lambda t}$

\textbf{Survival:} $P(T > t) = e^{-\lambda t}$

\textbf{Mean:} $E[T] = \dfrac{1}{\lambda}$ \quad (NOT $\lambda$!)

\textbf{Variance:} $\text{Var}(T) = \dfrac{1}{\lambda^2}$
\end{formulabox}

\begin{warningbox}
\textbf{Common Mistake:} Confusing $\lambda$ (rate) with $1/\lambda$ (mean)!

If rate is $\lambda = 2$ per minute, then mean time is $1/2 = 0.5$ minutes.
\end{warningbox}

\subsection{Memoryless Property}

\begin{criticalbox}
\textbf{The defining property of exponential distribution!}
\end{criticalbox}

\begin{formulabox}[colback=yellow!20,colframe=orange!80]
\textbf{MEMORYLESS PROPERTY:}
\[
\boxed{P(T > s + t \mid T > s) = P(T > t)}
\]

\textbf{In words:} If you've already waited $s$ time, the probability of waiting $t$ more is the same as from the start!

\textbf{Key fact:} Exponential is the ONLY continuous distribution with this property.
\end{formulabox}

\begin{examplebox}
\textbf{Problem:} Bus arrives exponentially with mean 10 minutes. You've waited 5 minutes.

What's the probability you wait 5 more minutes?

\textbf{Solution:}
\[
P(T > 10 \mid T > 5) = P(T > 5) = e^{-5/10} = e^{-0.5} \approx \boxed{0.606}
\]

The bus ``doesn't remember'' you've been waiting!
\end{examplebox}

\subsection{Minimum of Exponentials}

\begin{formulabox}
\textbf{MIN OF INDEPENDENT EXPONENTIALS:}

If $T_1 \sim \text{Exp}(\lambda_1)$ and $T_2 \sim \text{Exp}(\lambda_2)$ independent:

\vspace{0.2cm}

\textbf{Distribution of min:}
\[
\min(T_1, T_2) \sim \text{Exp}(\lambda_1 + \lambda_2)
\]

\textbf{Which happens first:}
\[
P(T_1 < T_2) = \frac{\lambda_1}{\lambda_1 + \lambda_2}
\]

\textbf{Expected time:}
\[
E[\min(T_1, T_2)] = \frac{1}{\lambda_1 + \lambda_2}
\]

\textbf{Connection:} This is identical to competing Poisson processes!
\end{formulabox}

\subsection{Maximum of Exponentials}

\begin{formulabox}
\textbf{MAX OF INDEPENDENT EXPONENTIALS:}

\vspace{0.2cm}

\textbf{Expected value:}
\[
E[\max(T_1, T_2)] = \frac{1}{\lambda_1} + \frac{1}{\lambda_2} - \frac{1}{\lambda_1 + \lambda_2}
\]

\textbf{Note:} No simple form for the full distribution.
\end{formulabox}

\begin{examplebox}
\textbf{Problem:} Bread making has two stages (happening simultaneously):
\begin{itemize}
  \item Rising: mean 8 hours
  \item Proofing: mean 4 hours
\end{itemize}

What's the expected time until both are done?

\textbf{Solution:} Need $E[\max(T_1, T_2)]$ where $\lambda_1 = 1/8$, $\lambda_2 = 1/4$:

\begin{align*}
E[\max] &= \frac{1}{1/8} + \frac{1}{1/4} - \frac{1}{1/8 + 1/4} \\
&= 8 + 4 - \frac{1}{3/8} \\
&= 12 - \frac{8}{3} = 12 - 2.67 \\
&= \boxed{9.33 \text{ hours}}
\end{align*}
\end{examplebox}

\subsection{Exponential Quick Reference}

\begin{formulabox}
\textbf{EXPONENTIAL DISTRIBUTION - FORMULA SHEET}

\vspace{0.2cm}

\textbf{Basic Facts ($T \sim \text{Exp}(\lambda)$):}
\begin{align*}
f(t) &= \lambda e^{-\lambda t} \\
P(T > t) &= e^{-\lambda t} \\
E[T] &= \frac{1}{\lambda}
\end{align*}

\textbf{Memoryless:} $P(T > s+t \mid T > s) = P(T > t)$

\textbf{Min:} $\min(T_1, T_2) \sim \text{Exp}(\lambda_1 + \lambda_2)$

\textbf{Max:} $E[\max(T_1, T_2)] = \dfrac{1}{\lambda_1} + \dfrac{1}{\lambda_2} - \dfrac{1}{\lambda_1 + \lambda_2}$
\end{formulabox}

\newpage

\section{BROWNIAN MOTION (Optional but Increasingly Common)}

\begin{criticalbox}
\textbf{Frequency:} Appears in \textbf{3/7 exams (43\%)}, but more common in recent years!
\end{criticalbox}

\subsection{What is Brownian Motion?}

\textbf{Continuous random movement:}
\begin{itemize}
  \item Particle jiggling in water
  \item Stock prices moving randomly
  \item Position of drunk walker
\end{itemize}

Like a random walk but in continuous time with infinitesimal steps.

\subsection{Standard Brownian Motion}

\begin{formulabox}
\textbf{STANDARD BROWNIAN MOTION $B(t)$:}

\vspace{0.2cm}

\textbf{Definition:}
\begin{enumerate}
  \item $B(0) = 0$ (starts at origin)
  \item $B(t) \sim N(0, t)$ (normally distributed)
  \item Independent increments
  \item $B(t) - B(s) \sim N(0, t-s)$ for $t > s$
\end{enumerate}

\textbf{Key Properties:}
\begin{align*}
E[B(t)] &= 0 \\
\text{Var}(B(t)) &= t \\
E[B(t)^2] &= t
\end{align*}

\textbf{Conditional:}
\begin{align*}
E[B(t) \mid B(s) = x] &= x \quad (t > s) \\
E[B(t)^2 \mid B(s) = x] &= x^2 + (t-s) \quad (t > s)
\end{align*}
\end{formulabox}

\subsection{Geometric Brownian Motion (Stock Prices)}

\begin{formulabox}
\textbf{GEOMETRIC BROWNIAN MOTION:}
\[
S(t) = S_0 e^{B(t)}
\]

Always positive (since $e^x > 0$), used to model stock prices.
\end{formulabox}

\begin{examplebox}
\textbf{Problem:} Stock price $S(t) = e^{B(t)}$ starting at $S(0) = 1$.

What's $P(S(2) \geq S(1)^2)$?

\textbf{Solution:}
\begin{align*}
S(2) \geq S(1)^2 &\Leftrightarrow e^{B(2)} \geq (e^{B(1)})^2 \\
&\Leftrightarrow e^{B(2)} \geq e^{2B(1)} \\
&\Leftrightarrow B(2) \geq 2B(1) \\
&\Leftrightarrow B(2) - 2B(1) \geq 0
\end{align*}

Now find distribution of $B(2) - 2B(1)$:
\begin{align*}
B(2) - 2B(1) &= [B(2) - B(1)] - B(1)
\end{align*}

where $B(2) - B(1) \sim N(0,1)$ and $B(1) \sim N(0,1)$ are independent.

Therefore $B(2) - 2B(1) \sim N(0, 1+1) = N(0, 2)$.

\begin{align*}
P(B(2) - 2B(1) \geq 0) &= P\left(\frac{B(2) - 2B(1)}{\sqrt{2}} \geq 0\right) \\
&= P(Z \geq 0) \quad \text{where } Z \sim N(0,1) \\
&= \boxed{0.5}
\end{align*}
\end{examplebox}

\newpage

\section{PATTERN RECOGNITION GUIDE}

\begin{criticalbox}
\textbf{The Secret to Success:} Recognize the problem type instantly, apply the right technique!
\end{criticalbox}

\subsection{When You See... Use This!}

\begin{table}[h]
\centering
\small
\begin{tabularx}{\textwidth}{lX}
\toprule
\textbf{Keywords/Setup} & \textbf{Technique to Use} \\
\midrule
``Poisson process'' + ``types/categories'' & \textbf{THINNING}: Type A rate = $\lambda p$ \\
\midrule
``which happens first'' + rates & \textbf{COMPETING PROCESSES}: $P(A \text{ first}) = \lambda_A/(\lambda_A + \lambda_B)$ \\
\midrule
``given that $n$ events occurred'' & \textbf{CONDITIONAL EXPECTATION}: $E[T|N=n] = n/\lambda$ \\
\midrule
``random walk'' + ``probability of reaching'' & \textbf{FIRST-STEP ANALYSIS} for hitting probability $p(i)$ \\
\midrule
``random walk'' + ``expected time/steps'' & \textbf{FIRST-STEP ANALYSIS} for expected time $\mu(i)$ (DON'T FORGET +1!) \\
\midrule
``joint density'' or ``$f(x,y)$'' & Set up \textbf{JOINT PDF}, use $f_{X,Y} = f_{Y|X} \cdot f_X$ \\
\midrule
``given $X=x$, find $E[Y]$'' & Use \textbf{LAW OF TOTAL EXPECTATION}: $E[Y] = E[E[Y|X]]$ \\
\midrule
``$Y = UX$'' where $U \sim$ Unif$[0,1]$ & $f_{Y|X}(y|x) = 1/x$ for $0 < y < x$ \\
\midrule
``exponential'' + ``min'' or ``first'' & \textbf{MIN OF EXPONENTIALS}: $\min \sim \text{Exp}(\lambda_1 + \lambda_2)$ \\
\midrule
``exponential'' + ``max'' or ``both done'' & $E[\max] = 1/\lambda_1 + 1/\lambda_2 - 1/(\lambda_1+\lambda_2)$ \\
\midrule
``Brownian motion'' or ``$B(t)$'' & Use BM properties: $E[B(t)] = 0$, $\text{Var}(B(t)) = t$ \\
\midrule
``$S(t) = e^{B(t)}$'' (stock prices) & \textbf{GEOMETRIC BM}: Convert to statements about $B(t)$ \\
\bottomrule
\end{tabularx}
\end{table}

\subsection{Problem Type Flowchart}

\begin{enumerate}[leftmargin=*]
  \item \textbf{Does it mention ``arrivals'' or ``rate''?}
    \begin{itemize}
      \item YES → Probably Poisson process
      \item Check for types/categories → Use thinning
      \item Check for ``which first'' → Competing processes
    \end{itemize}

  \item \textbf{Does it mention ``positions'', ``game'', ``walk''?}
    \begin{itemize}
      \item YES → Random walk with first-step analysis
      \item Asked for probability? → Set up $p(i)$ equations
      \item Asked for expected time? → Set up $\mu(i)$ equations
    \end{itemize}

  \item \textbf{Does it involve TWO random variables?}
    \begin{itemize}
      \item YES → Joint/conditional densities
      \item Check for $Y = UX$ → Use standard formulas
      \item Need expectation? → Try law of total expectation
    \end{itemize}

  \item \textbf{Does it mention ``waiting time'' or ``lifetime''?}
    \begin{itemize}
      \item YES → Probably exponential
      \item Multiple waiting times? → Min/max of exponentials
    \end{itemize}

  \item \textbf{Does it mention ``Brownian motion'' or $B(t)$?}
    \begin{itemize}
      \item YES → Use BM properties
      \item Stock prices? → Geometric BM
    \end{itemize}
\end{enumerate}

\newpage

\section{COMPLETE FORMULA SHEET}

\subsection{Poisson Processes}

\begin{tcolorbox}[colback=formulacolor!5,colframe=formulacolor]
\textbf{Basic:} $N(t) \sim \text{Poisson}(\lambda t)$, \quad $E[N(t)] = \text{Var}(N(t)) = \lambda t$

\[
P(N(t) = k) = \frac{(\lambda t)^k e^{-\lambda t}}{k!}
\]

\textbf{Thinning:} Type A prob $p$ → Type A is Poisson($\lambda p$)

\textbf{Competing:} $P(A \text{ before } B) = \dfrac{\lambda_A}{\lambda_A + \lambda_B}$, \quad $E[T] = \dfrac{1}{\lambda_A + \lambda_B}$

\textbf{Conditional:} $E[T \mid N(T) = n] = \dfrac{n}{\lambda}$
\end{tcolorbox}

\subsection{Random Walks}

\begin{tcolorbox}[colback=formulacolor!5,colframe=formulacolor]
\textbf{Hitting Probability:}
\[
p(i) = p_{\uparrow} p(i+1) + p_{\downarrow} p(i-1) + p_{\leftrightarrow} p(i)
\]
Boundaries: $p(a) = 0$, $p(b) = 1$

\textbf{Expected Time:}
\[
\mu(i) = \mathbf{1} + p_{\uparrow} \mu(i+1) + p_{\downarrow} \mu(i-1) + p_{\leftrightarrow} \mu(i)
\]
Boundaries: $\mu(a) = 0$, $\mu(b) = 0$

\textbf{Symmetric Shortcuts ($p_{\uparrow} = p_{\downarrow} = 1/2$):}
\[
p(i) = \frac{i-a}{b-a}, \quad \mu(i) = (i-a)(b-i)
\]
\end{tcolorbox}

\subsection{Joint \& Conditional Densities}

\begin{tcolorbox}[colback=formulacolor!5,colframe=formulacolor]
\[
f_{X|Y}(x|y) = \frac{f_{X,Y}(x,y)}{f_Y(y)}, \quad f_{X,Y}(x,y) = f_{Y|X}(y|x) \cdot f_X(x)
\]
\[
f_X(x) = \int f_{X,Y}(x,y) \, dy
\]

\textbf{Law of Total Expectation:} $E[Y] = E[E[Y|X]]$

\textbf{For $Y = UX$, $U \sim$ Unif$[0,1]$:}
\[
f_{Y|X}(y|x) = \frac{1}{x} \text{ for } 0 < y < x, \quad E[Y|X=x] = \frac{x}{2}
\]
\end{tcolorbox}

\subsection{Exponential Distribution}

\begin{tcolorbox}[colback=formulacolor!5,colframe=formulacolor]
\textbf{$T \sim \text{Exp}(\lambda)$:} $f(t) = \lambda e^{-\lambda t}$, \quad $P(T > t) = e^{-\lambda t}$, \quad $E[T] = \dfrac{1}{\lambda}$

\textbf{Memoryless:} $P(T > s+t \mid T > s) = P(T > t)$

\textbf{Min:} $\min(T_1, T_2) \sim \text{Exp}(\lambda_1 + \lambda_2)$, \quad $P(T_1 < T_2) = \dfrac{\lambda_1}{\lambda_1 + \lambda_2}$

\textbf{Max:} $E[\max(T_1, T_2)] = \dfrac{1}{\lambda_1} + \dfrac{1}{\lambda_2} - \dfrac{1}{\lambda_1 + \lambda_2}$
\end{tcolorbox}

\subsection{Brownian Motion}

\begin{tcolorbox}[colback=formulacolor!5,colframe=formulacolor]
\textbf{$B(t)$ standard BM:} $E[B(t)] = 0$, \quad $\text{Var}(B(t)) = t$, \quad $B(t) \sim N(0,t)$

$B(t) - B(s) \sim N(0, t-s)$ for $t > s$

\textbf{Conditional:} $E[B(t)^2 \mid B(s) = x] = x^2 + (t-s)$ for $t > s$

\textbf{Geometric BM:} $S(t) = S_0 e^{B(t)}$ (stock prices)
\end{tcolorbox}

\newpage

\section{COMMON MISTAKES \& HOW TO AVOID THEM}

\subsection{Top 10 Mistakes Students Make}

\begin{enumerate}[leftmargin=*]
  \item \textbf{Forgetting ``+1'' in random walk expected time}
    \begin{itemize}
      \item Formula: $\mu(i) = \mathbf{1} + \cdots$ (THE 1 IS CRITICAL!)
      \item Fix: Write it in red on your cheat sheet
    \end{itemize}

  \item \textbf{Not using thinning when events split into types}
    \begin{itemize}
      \item Keyword: ``40\% order coffee, 30\% tea...''
      \item Fix: Always check if arrivals split by category
    \end{itemize}

  \item \textbf{Confusing $\lambda$ with $1/\lambda$ for exponential}
    \begin{itemize}
      \item Rate $\lambda = 2$/min means mean time $= 1/2 = 0.5$ min
      \item Fix: Write ``$E[T] = 1/\lambda$'' on cheat sheet
    \end{itemize}

  \item \textbf{Wrong boundary conditions for random walks}
    \begin{itemize}
      \item For $p(i)$: $p(a) = 0$, $p(b) = 1$
      \item For $\mu(i)$: $\mu(a) = 0$, $\mu(b) = 0$
      \item Fix: Always write boundaries first
    \end{itemize}

  \item \textbf{Not specifying range for joint PDFs}
    \begin{itemize}
      \item Must say ``for $0 < y < x$'' or similar
      \item Fix: Draw a picture of the region
    \end{itemize}

  \item \textbf{Not checking if random walk is symmetric}
    \begin{itemize}
      \item If $p_{\uparrow} = p_{\downarrow} = 1/2$, use shortcuts!
      \item Fix: First check symmetry before setting up equations
    \end{itemize}

  \item \textbf{Arithmetic errors in solving systems}
    \begin{itemize}
      \item First-step analysis needs solving 2-3 equations
      \item Fix: Show all work, check answer makes sense
    \end{itemize}

  \item \textbf{Using $\lambda t$ instead of $\lambda$ in Poisson}
    \begin{itemize}
      \item Rate is $\lambda$, expected count in time $t$ is $\lambda t$
      \item Fix: Be clear about units (per minute vs total)
    \end{itemize}

  \item \textbf{Forgetting independence assumptions}
    \begin{itemize}
      \item Thinning requires independence
      \item Min of exponentials requires independence
      \item Fix: Check problem statement
    \end{itemize}

  \item \textbf{Not showing work / skipping steps}
    \begin{itemize}
      \item Partial credit is huge!
      \item Fix: Write every step, even if obvious
    \end{itemize}
\end{enumerate}

\newpage

\section{EXAM DAY STRATEGY}

\subsection{Before the Exam}

\begin{enumerate}[leftmargin=*]
  \item \textbf{Sleep well} - Seriously, don't pull an all-nighter
  \item \textbf{Prepare cheat sheet} - One page, both sides, handwritten
  \item \textbf{Eat something} - Low energy = mistakes
  \item \textbf{Arrive early} - Get settled, calm down
\end{enumerate}

\subsection{During the Exam}

\begin{strategybox}
\textbf{FIRST 5 MINUTES:}
\begin{enumerate}
  \item Skim all questions quickly
  \item Identify problem types (Poisson? Random walk? Joint density?)
  \item Star the ones you feel confident about
  \item Budget time: 35 minutes per question
\end{enumerate}
\end{strategybox}

\begin{strategybox}
\textbf{FOR EACH PROBLEM:}
\begin{enumerate}
  \item \textbf{Read carefully} - What's given? What's asked?
  \item \textbf{Identify technique} - Use pattern recognition guide
  \item \textbf{Write what you know} - Definitions, parameters, distributions
  \item \textbf{Show all work} - Partial credit!
  \item \textbf{Check units} - Per minute? Per hour? Total?
  \item \textbf{Simplify answer} - Reduce fractions
  \item \textbf{Sanity check} - Does it make sense?
\end{enumerate}
\end{strategybox}

\subsection{Time Management}

\begin{itemize}[leftmargin=*]
  \item \textbf{Don't get stuck} - If blocked after 5 min, move on
  \item \textbf{Do easy parts first} - Get partial credit
  \item \textbf{Come back to hard parts} - Fresh perspective helps
  \item \textbf{Leave 5 minutes} - Final check, upload
\end{itemize}

\subsection{Partial Credit Strategy}

Even if you can't solve fully:
\begin{itemize}[leftmargin=*]
  \item Write down relevant formulas
  \item Set up the problem (define variables)
  \item Explain your approach in words
  \item Show partial calculations
  \item Write ``If I could solve X, then I would do Y''
\end{itemize}

\newpage

\section{STUDY PLAN}

\subsection{4-Day Intensive Plan}

\subsubsection{Day 1: Poisson Processes (8 hours)}

\textbf{Morning (4 hours):}
\begin{itemize}
  \item Read Section 1 of this guide (2 hours)
  \item Work through all examples yourself (2 hours)
\end{itemize}

\textbf{Afternoon (4 hours):}
\begin{itemize}
  \item Practice: F2025 Q3 (Coffee shop)
  \item Practice: S2025 Q3 (Taxi stop)
  \item Practice: 2019 Q1 (Plagues)
  \item Check solutions carefully
\end{itemize}

\textbf{Evening Check:}
\begin{itemize}
  \item Can you explain thinning without notes?
  \item Can you solve competing process problems?
\end{itemize}

\subsubsection{Day 2: Random Walks (8 hours)}

\textbf{Morning (4 hours):}
\begin{itemize}
  \item Read Section 2 of this guide (2 hours)
  \item Practice writing first-step equations (2 hours)
\end{itemize}

\textbf{Afternoon (4 hours):}
\begin{itemize}
  \item Practice: F2025 Q2 (Sleepy gambler)
  \item Practice: S2025 Q1 (Flight prices)
  \item Practice: 2017 Q3 (Asymmetric game)
\end{itemize}

\textbf{Evening Check:}
\begin{itemize}
  \item Write formulas from memory
  \item Do you remember the ``+1''?
\end{itemize}

\subsubsection{Day 3: Joint Densities + Full Exam (8 hours)}

\textbf{Morning (3 hours):}
\begin{itemize}
  \item Read Section 3 (2 hours)
  \item Practice: F2025 Q1 (Insurance) (1 hour)
\end{itemize}

\textbf{Afternoon (5 hours):}
\begin{itemize}
  \item Full timed practice exam: F2025 or S2025 (2 hours)
  \item Check solutions thoroughly (2 hours)
  \item Quick review of exponential \& BM (1 hour)
\end{itemize}

\subsubsection{Day 4: Review + Weak Spots (6 hours)}

\textbf{Morning (3 hours):}
\begin{itemize}
  \item Another full practice exam
  \item Identify your 2-3 weakest topics
\end{itemize}

\textbf{Afternoon (3 hours):}
\begin{itemize}
  \item Redo problems you got wrong
  \item Focus on weak topics only
  \item Make cheat sheet
\end{itemize}

\textbf{Evening:}
\begin{itemize}
  \item Light review of formula sheet
  \item Get good sleep!
\end{itemize}

\subsection{2-Day Crash Plan}

If you only have 2 days:

\subsubsection{Day 1: Poisson + Random Walks (12 hours)}
\begin{itemize}
  \item These appear in 100\% of exams
  \item Read Sections 1 \& 2
  \item Do 3-4 practice problems each
  \item Master thinning and first-step analysis
\end{itemize}

\subsubsection{Day 2: Everything Else + Practice (12 hours)}
\begin{itemize}
  \item Skim Section 3 (joint densities)
  \item Do 2 full practice exams
  \item Focus on understanding solutions
\end{itemize}

\textbf{Skip:} Brownian motion if pressed for time

\newpage

\section{FINAL THOUGHTS}

\subsection{You Can Do This!}

Remember:
\begin{itemize}[leftmargin=*]
  \item This exam is \textbf{predictable} - same topics every year
  \item Master \textbf{3 techniques} = handle 75\% of exam
  \item Understanding \textgreater\ memorizing
  \item Partial credit is huge
  \item Thousands have learned this before you
\end{itemize}

\subsection{The Formula for Success}

\begin{tcolorbox}[colback=yellow!20,colframe=orange!80,arc=3mm]
\centering
\Large
\textbf{Success = Pattern Recognition + Practice + Showing Work}

\vspace{0.3cm}

\normalsize
\begin{itemize}[leftmargin=2cm]
  \item \textbf{Pattern Recognition:} See ``types'' → Think thinning
  \item \textbf{Practice:} Do problems, check solutions, understand mistakes
  \item \textbf{Show Work:} Write every step = maximum partial credit
\end{itemize}
\end{tcolorbox}

\subsection{What to Focus On}

\begin{table}[h]
\centering
\begin{tabularx}{0.9\textwidth}{Xc}
\toprule
\textbf{Topic} & \textbf{Time Investment} \\
\midrule
Poisson Processes (esp. thinning) & 30\% \\
Random Walks (first-step analysis) & 30\% \\
Joint/Conditional Densities & 20\% \\
Exponential Distribution & 10\% \\
Practice Exams & 10\% \\
\bottomrule
\end{tabularx}
\end{table}

\subsection{Last-Minute Checklist}

\textbf{Before Exam:}
\begin{itemize}[label=$\square$]
  \item Cheat sheet prepared (formulas, NOT examples)
  \item Calculator ready (if allowed)
  \item Reviewed formula sheet (Section 8)
  \item Did at least 2 full practice exams
  \item Got good sleep
\end{itemize}

\textbf{During Exam:}
\begin{itemize}[label=$\square$]
  \item Read all questions first
  \item Budget 35 minutes per question
  \item Show all work (partial credit!)
  \item Check answers make sense
  \item Upload with 5 minutes to spare
\end{itemize}

\vfill

\begin{center}
\Large\textbf{Good luck!}

\vspace{0.5cm}

\normalsize
You've got the tools. You've got the guide. Now go ace this exam!

\vspace{1cm}

\textit{``Success is not final, failure is not fatal: it is the courage to continue that counts.''}

--- Winston Churchill
\end{center}

\end{document}
